\begin{minipage}{0.65\linewidth}
    On fixe le solide $\mathrm{(S)}$ précédent à un ressort horizontal à spires
    non jointives, de masse négligeable et de constante de raideur $K$.
    
    À l'équilibre, le centre d'inertie $G$ coïncide avec l'origine du repère
    $(O,\ex)$ lié à la terre considéré comme galiléen
    (Fig.~\ref{fig:solide_ressort}).
    
    On écarte le solide $\mathrm{(S)}$ de sa position d'équilibre et on le
    libère sans vitesse initiale à l'instant $t_{0}=0$.
\end{minipage}
\hfill
\begin{minipage}{0.33\linewidth}
    \begin{figure}[H]
        \centering
        \begin{tikzpicture}
    \node[inner sep=0pt,minimum width=4.5cm,minimum height=3mm,
        top color=black!70] (sol) {};
    % Axe xx'
    \coordinate (O) at ($(sol.south west)!.8!(sol.south east)$);
    \draw[-Latex] (sol.south west)
        node[below right,inner xsep=0pt] {$x^{\prime}$} -- (O)
        node[below=0.4mm] {$O$} -- (sol.south east) -- ++(0.7,0)
        node[below left,inner xsep=0pt,inner ysep=2mm] {$x$};
    % Fixed
    \node[inner sep=0pt,draw,minimum width=3mm,minimum height=7mm,
        anchor=south,pattern=north west lines,
        thick] (fix) at ([xshift=3mm,yshift=0.4pt]sol.north west) {};
    % Solide (S)
    \node[inner sep=0pt,draw,minimum width=7mm,minimum height=7mm,
        anchor=south,rounded corners=2pt,thick,label=$\mathrm{(S)}$,
        fill=lightgray] (S) at
        ([yshift=0.4pt]O |- sol.north) {};
        %([xshift=-6mm,yshift=0.4pt]sol.north east) {};
    \node[inner sep=1pt,circle,fill,
        label={[right=-0.8mm,font=\scriptsize]:$G$}] (G) at (S.center) {};
    \draw[{Rectangle[length=1pt,width=4pt]}-{Stealth[length=5pt]},
        thick] (O) -- ++(0.8,0)
        node[below left,inner xsep=0pt] {$\ex$};
    % Ressort
    \draw[decoration={
          coil,
          aspect=0.3,
          segment length=2.0mm,
          amplitude=2.4mm,
          pre length=1mm,
          post length=1mm},
          decorate,thick,black] (fix.east) -- (S.west);
\end{tikzpicture}

        \caption{}
        \label{fig:solide_ressort}
    \end{figure}
\end{minipage}

\begin{donnees}
    \begin{itemize}
      \item Tous les frottements sont négligeables~;
      \item On choisit l'état où le ressort n'est pas déformé comme référence de
            l'énergie potentielle élastique $E_{pe}$ et le plan horizontal
            contenant $G$ comme état de référence de l'énergie potentielle de
            pesanteur $E_{pp}$.
    \end{itemize}
\end{donnees}
\begin{minipage}[t][][t]{0.68\linewidth}
    La courbe de la Fig.~\ref{fig:2bac_svt_national_2016_normale_phys_ex3_3}
    représente les variations de $E_{pe}$ en fonction de $x^{2}$, carré de
    l'abscisse $x$ du centre d'inertie $G$ dans le repère $(O,\ex)$.
    \begin{enumerate}
        \item En exploitant la courbe de la figure~(\ref{fig:Epe_de_x2}),
              trouver les valeurs de~:
              \begin{enumerate}
                  \item la constante de raideur $K$.
                  \item l'énergie potentielle élastique maximale $E_{pe,\max}$.
                  \item l'amplitude $X_{m}$ des oscillations.
              \end{enumerate}
        \item Déduire, en justifiant votre réponse, la valeur de l'énergie
              mécanique $E_{m}$ du système oscillant.
    \end{enumerate}
\end{minipage}
\hfill
\begin{minipage}[t][][t]{0.31\linewidth}
    \begin{figure}[H]
        \centering
        \begin{tikzpicture}
            \def\K{1}
            \begin{axis}[
                    xmin=0,xmax=19,
                    ymin=0,ymax=9.8,
                    width=6.5cm,height=6.5cm,
                    xlabel=$x^{2}~(\times\qty{1e-4}{\square\meter})$,
                    ylabel=$E_{pe}~(\unit{\mJ})$,
                    scaled ticks=false,
                    xtick distance=4, ytick distance=2,
                    xticklabels={,,4}, yticklabels={,,2},
                ]
                \addplot[domain=0:16,samples=200,
                    very thick,blue] {0.5*\K*x}
                    coordinate[pos=1](A);
            \end{axis}
            \draw[very thick,blue] ([yshift=-5mm]A) -- ([yshift=5mm]A);
            \foreach \f in {0.02,0.14,...,1}
                \draw[very thick,blue]
                    ($([yshift=-5mm]A)!\f!([yshift=5mm]A)$) -- ++(20:.2);
        \end{tikzpicture}
        \caption{}
        \label{fig:2bac_svt_national_2016_normale_phys_ex3_3}
    \end{figure}
\end{minipage}
\begin{enumerate}
    \setcounter{enumi}{2}
    \item Le centre d'inertie $G$ passe par la position d'équilibre dans le sens
          positif avec la vitesse $v=\qty{0,25}{\v}$.

          Montrer que l'expression de la période propre des oscillations s'écrit
          $T_{0}=2\pi\cdot\frac{X_{m}}{v}$. Calculer $T_{0}$.
\end{enumerate}

